\documentclass[11pt,a4paper]{article}
\usepackage[margin=19mm,top=17mm,bottom=17mm]{geometry}
\usepackage{amsmath,amssymb,amsfonts}
\usepackage{mathptmx}
\usepackage{enumitem}
\usepackage{fancyhdr}
\usepackage{draftwatermark}
\usepackage{graphicx}
\usepackage{array,tabularx,booktabs}
\usepackage[hidelinks]{hyperref}
\setlength{\parindent}{0pt}\setlength{\parskip}{4pt}
\setlist[enumerate,1]{leftmargin=1.1em,itemsep=3pt,parsep=2pt}
\setlist[itemize]{leftmargin=1.4em,itemsep=1pt}
\pagestyle{fancy}\fancyhf{}
\fancyhead[L]{\footnotesize \textbf{Marks : 150}}
\fancyhead[C]{\footnotesize \textbf{TEST ON DEFINITE INTEGRALS}}
\fancyhead[R]{\footnotesize \textbf{Time : 2 : 30 hours}}
\fancyfoot[L]{\footnotesize Think C}
\fancyfoot[C]{\footnotesize Education is not just about Information  $\cdot$  Page \thepage}
\fancyfoot[R]{\footnotesize }
\renewcommand{\headrulewidth}{0.4pt}
\SetWatermarkText{ThinkC}
\SetWatermarkScale{1.4}
\SetWatermarkAngle{45}
\SetWatermarkLightness{82}
\begin{document}
\vspace{2mm}
\noindent \textbf{Marks : 60} \hfill \textbf{Time : 1 : 00 hours}
\noindent\rule{\linewidth}{1.1pt}\vspace{-1mm}\noindent\rule{\linewidth}{0.5pt}
{\centering \Large \textbf{STRAIGHT LINES PAPER 2}\par}
\noindent\rule{\linewidth}{1.1pt}\vspace{-1mm}\noindent\rule{\linewidth}{0.5pt}
\vspace{2mm}

\begin{enumerate}[resume]
\item In the triangle with vertices \((0,0)\), \((8,6)\), and \((20,0)\), a square is placed with its lower side along the \(x\)-axis and its two upper vertices lying on the two sloping sides of the triangle. The area of the square is \hfill [4 marks]
{\renewcommand{\arraystretch}{1.8}
\begin{tabularx}{\linewidth}{X X X X}
(A) \quad \(\frac{3600}{169}\) & (B) \quad \(\frac{400}{13}\) & (C) \quad \(\frac{900}{49}\) & (D) \quad \(\frac{1600}{121}\) \\[10pt]
\end{tabularx}}
\item Let \(A(-2,1)\) and \(B(4,5)\) be fixed points in the Cartesian plane. If a point \(C(x,y)\) satisfies \(AC=2BC\), then the number of distinct non-degenerate triangles \(ABC\) that can be formed is: \hfill [1 marks]
{\renewcommand{\arraystretch}{1.8}
\begin{tabularx}{\linewidth}{X X X X}
(A) \quad Zero & (B) \quad One & (C) \quad Two & (D) \quad Infinitely many \\[10pt]
\end{tabularx}}
\item A light ray starts from the point \(A(-2,3)\), strikes the \(x\)-axis at \(B\), and after reflection travels through the point \(P(4,5)\). The equation of the incident ray \(AB\) is: \hfill [4 marks]
{\renewcommand{\arraystretch}{1.8}
\begin{tabularx}{\linewidth}{X X X X}
(A) \quad \(4x+3y=1\) & (B) \quad \(3x+4y=6\) & (C) \quad \(4x-3y=-1\) & (D) \quad \(3x-4y=-6\) \\[10pt]
\end{tabularx}}
\item The vertices of a parallelogram \(ABCD\) are \(A(-2,1)\), \(B(6,4)\), \(C(10,12)\), and \(D(2,9)\). A straight line passing through the origin divides the parallelogram into two congruent regions. Which of the following can be the equation of this line? \hfill [4 marks]
{\renewcommand{\arraystretch}{1.8}
\begin{tabularx}{\linewidth}{X X X X}
(A) \quad \(13x-8y=0\) & (B) \quad \(8x+13y=0\) & (C) \quad \(13x+8y=0\) & (D) \quad \(8x-13y=0\) \\[10pt]
\end{tabularx}}
\item The fixed vertices of a triangle \(PQR\) are \(P(1,4)\) and \(Q(-2,7)\), while the vertex \(R\) moves on the line \(5x-12y+9=0\). If \(G\) is the centroid of \(\triangle PQR\), then the locus of \(G\) is: \hfill [4 marks]
{\renewcommand{\arraystretch}{1.8}
\begin{tabularx}{\linewidth}{X X X X}
(A) \quad \(15x-36y+146=0\) & (B) \quad \(15x-36y-146=0\) & (C) \quad \(5x-12y+146=0\) & (D) \quad \(5x-12y-146=0\) \\[10pt]
\end{tabularx}}
\item In a triangle \(ABC\), the point \(A\) is \((3,4)\). The equations of the altitude and the internal angle bisector drawn from \(B\) are respectively \(4x+y-5=0\) and \(x-y=0\). The equation of the side \(BC\) is: \hfill [1 marks]
{\renewcommand{\arraystretch}{1.8}
\begin{tabularx}{\linewidth}{X X X X}
(A) \quad \(2x-3y+1=0\) & (B) \quad \(3x-2y-1=0\) & (C) \quad \(x+4y-5=0\) & (D) \quad \(3x+2y-5=0\) \\[10pt]
\end{tabularx}}
\item The pair of straight lines represented by \(2x^2+6xy+5y^2=0\) is reflected in the line \(y=x\). The equation of the image pair is: \hfill [4 marks]
{\renewcommand{\arraystretch}{1.8}
\begin{tabularx}{\linewidth}{X X X X}
(A) \quad \(5x^2+6xy+2y^2=0\) & (B) \quad \(2x^2-6xy+5y^2=0\) & (C) \quad \(5x^2-6xy+2y^2=0\) & (D) \quad \(2x^2+6xy+5y^2=0\) \\[10pt]
\end{tabularx}}
\item For all real numbers \(p,q\), not both zero, the line \((2p-q)x+(p+3q)y=5p+7q\) belongs to a family of lines passing through a fixed point. The coordinates of this fixed point are: \hfill [4 marks]
{\renewcommand{\arraystretch}{1.8}
\begin{tabularx}{\linewidth}{X X}
(A) \quad \(\left(\frac{8}{7},\frac{19}{7}\right)\) & (B) \quad \(\left(\frac{19}{7},\frac{8}{7}\right)\) \\[10pt]
(C) \quad \(\left(\frac{8}{5},\frac{19}{5}\right)\) & (D) \quad \(\left(\frac{5}{7},\frac{7}{5}\right)\) \\[10pt]
\end{tabularx}}
\item A triangle has vertices \((-3,2)\), \((5,10)\), and \((13,2)\). A square is inscribed in the triangle such that one of its sides lies on the line \(y=2\), while the two opposite vertices lie on the other two sides of the triangle. The area of the square is: \hfill [4 marks]
{\renewcommand{\arraystretch}{1.8}
\begin{tabularx}{\linewidth}{X X X X}
(A) \quad \(\dfrac{256}{9}\) & (B) \quad \(\dfrac{64}{9}\) & (C) \quad \(\dfrac{512}{9}\) & (D) \quad \(\dfrac{256}{7}\) \\[10pt]
\end{tabularx}}
\item Let \(P(-2,1)\) and \(Q(4,-3)\) be fixed points in the Cartesian plane. A point \(R\) is chosen such that \(PR=2QR\). Considering only non-degenerate triangles \(PQR\), which statement is correct? \hfill [1 marks]
{\renewcommand{\arraystretch}{1.8}
\begin{tabularx}{\linewidth}{X X}
(A) \quad Exactly one such triangle is possible. & (B) \quad Exactly two such triangles are possible. \\[10pt]
(C) \quad No such triangle is possible. & (D) \quad Infinitely many such triangles are possible. \\[10pt]
\end{tabularx}}
\item A ray of light passes through the point \(A(-2,3)\), strikes the \(x\)-axis at \(B\), and after reflection passes through the point \(C(4,5)\). The equation of the incident ray \(AB\) is: \hfill [4 marks]
{\renewcommand{\arraystretch}{1.8}
\begin{tabularx}{\linewidth}{X X X X}
(A) \quad \(4x+3y=1\) & (B) \quad \(3x+4y=6\) & (C) \quad \(4x-3y=-17\) & (D) \quad \(3x-4y=-18\) \\[10pt]
\end{tabularx}}
\item The vertices of a parallelogram \(ABCD\) are \(A(4,2)\), \(B(12,5)\), \(C(15,13)\) and \(D(7,10)\). A line passing through the origin divides the parallelogram into two congruent parts. The slope of this line is \hfill [1 marks]
{\renewcommand{\arraystretch}{1.8}
\begin{tabularx}{\linewidth}{X X X X}
(A) \quad \(\frac{15}{19}\) & (B) \quad \(\frac{19}{15}\) & (C) \quad \(\frac{5}{8}\) & (D) \quad \(\frac{8}{5}\) \\[10pt]
\end{tabularx}}
\item Let the fixed vertices of a triangle be \(A(1,4)\) and \(B(7,-2)\), while the third vertex \(C\) moves on the line \(5x-12y+9=0\). If \(G(X,Y)\) is the centroid of \(\triangle ABC\), then the locus of \(G\) is: \hfill [4 marks]
{\renewcommand{\arraystretch}{1.8}
\begin{tabularx}{\linewidth}{X X X X}
(A) \quad \(15X-36Y-7=0\) & (B) \quad \(15X-36Y+7=0\) & (C) \quad \(5X-12Y-7=0\) & (D) \quad \(5X-12Y+9=0\) \\[10pt]
\end{tabularx}}
\item In a non-degenerate triangle \(ABC\), the point \(A\) is \((1,4)\). The altitude and the internal angle bisector drawn from \(B\) have equations \(x+y-5=0\) and \(x-2y+1=0\), respectively. The equation of the side \(BC\) is \hfill [4 marks]
{\renewcommand{\arraystretch}{1.8}
\begin{tabularx}{\linewidth}{X X X X}
(A) \quad \(7x+y-23=0\) & (B) \quad \(7x-y-19=0\) & (C) \quad \(x-7y+11=0\) & (D) \quad \(3x+2y-13=0\) \\[10pt]
\end{tabularx}}
\end{enumerate}
\end{document}